\section{Introduction}
\label{sec:intro}

How much of the arrow of time does an observer lose by looking at a system through a blur?
We ask this for the simplest unstable motion, the pendulum near its upper equilibrium, $\dot z=p$, $\dot p=\sin z$, which pulls nearby states apart at the rate $\kappa=1$.
The observer does not see the state $X$ itself, only $Y=X+\sigma Z$ with $Z\sim N(0,I_2)$: a Gaussian blur of width $\sigma$.
From the blurred density the observer reconstructs a drift through its score, and the rate $R$ that this drift reaches is smaller than $\kappa$.
The resolution gap $g=(\kappa-R)/\sigma$ measures the loss per unit of resolution (Sec.~\ref{sec:main} gives the precise definition).
The question is how small $g$ can become as the blur gets finer, and whether the answer depends on the state.

The natural candidate for a uniform answer is the logarithmic law
\begin{equation}
g\;\ge\;\frac{c'}{\sqrt{\ell+C'}},\qquad \ell=\ln(1/\sigma),
\label{eq:conj1}
\end{equation}
with constants $c',C'>0$ that do not depend on the state.
It was conjectured in Ref.~\cite{E}, where the weaker bound $g\ge c/\ell$ was proved for states without fine structure.
The law allows the gap to shrink as the blur gets finer, but only logarithmically, whatever the state.

So far the lower side of~\eqref{eq:conj1} was known only for states built from windows~\cite{F,H,I,J}.
A window is a smooth envelope that carries a periodic ripple in the stable direction of the equilibrium; the results cover weak and strong ripples and every place where the observer focuses.
A state that is not a window, for example one whose information sits in many scattered spots, was outside them.

This paper removes the shape.
Theorem~\ref{thm:main} gives a lower bound on $g$ for every probability measure on the plane with a finite first moment that is confined to the neighbourhood of the equilibrium, $\Var(X_z)\le\pi^2$.
The bound decays like $\ell^{-5/4}(\ln\ell)^{3/4}$.
That is slower than any power of $\sigma$, so the gap never closes at a finite resolution, but faster than the conjectured $\ell^{-1/2}$.
It holds for $\sigma\le1.5\times10^{-14}$, and its constant is small.
We regard the theorem as the first unconditional step towards~\eqref{eq:conj1}, and its power as the present one.
Section~\ref{sec:outlook} shows which inequality each of the next steps rests on: sharpening the power to $\ell^{-1/2}$, and widening the range of $\sigma$.

\paragraph*{Related work.}
That the arrow of time appears on the side of the coarse-grained description is a classical view~\cite{Lebowitz}.
The present problem is one concrete instance of it: the coarse-graining is a Gaussian blur of width $\sigma$, and the question is how much of the arrow it leaves, uniformly over the states.
A Gaussian blur moves the density along the heat flow, and how Fisher information behaves along that flow is the subject of de~Bruijn's identity and Stam's inequality~\cite{Stam}; for Gaussian channels, mutual information and the minimum mean-square error are related as well~\cite{GSV}.
The quantity $F(t)$ of Sec.~\ref{sec:circles}, the information in the stable direction along a heat flow in the unstable one, belongs to this family.
Tweedie's formula~\cite{Robbins,Efron}, which connects the posterior mean with the score, also underlies score-based estimation~\cite{Hyvarinen} and diffusion models~\cite{SongSDE}.
The central object of this paper, $\st=-\E[Z\mid Y]$, is exactly that score.
We use these connections as background only; the results of this paper concern the pendulum and the gap.

\paragraph*{What is new.}
Three steps make the shape irrelevant.
\begin{enumerate}
\item A bound on the effective wavenumber, $\keff^2\le\psi(\ln(1/A))$ with $\psi(L)=289L^{5/2}/(\ln2L)^{3/2}+4L+3.57$, for every state (Theorem~\ref{thm:circles}).
Earlier proofs needed such a bound and had it only as a conjecture or for windows.
The proof goes through complex time: the stable-direction score of the density with $v$-noise $z$ is a holomorphic function of $z$ with values in a Hilbert space, and Hadamard's three-circle theorem transports its smallness at $z=1$ to $z$ slightly below one.
The three-circle argument was proposed by Google Gemini, one of the AI assistants used in this work (see the Acknowledgments), and checked here line by line and numerically at every step.
\item A threshold for the remainder of the gap formula of Ref.~\cite{E}: once the $w$-information and the curved $v$-information together exceed $112\sigma^2$, the remainder is at most half of the main terms (Lemma~\ref{lem:EF2}).
So every remainder estimate is needed only on an open set $\Oset$ of states with information of order $\sigma^2$.
\item On $\Oset$ the gap is expanded around the observed point $Y_z$ instead of the posterior mean $m_z$ (Lemma~\ref{lem:EGD}).
The state-independent remainder $(\sqrt{960}/6)\sigma^2$ of the old expansion disappears; the new remainder carries $\sqrt{A+B}$, and the single second-order quantity left, $\Wzz=\sigma^2\partial_z^2\varrho/\varrho$, is controlled by the same three circles in the isotropic flow (Lemmas~\ref{lem:iso} and~\ref{lem:sm}).
\end{enumerate}
The rest of the chain is borrowed from Refs.~\cite{E,F} and from the lab stages behind this paper; Section~\ref{sec:chain} shows how it fits together, and the appendices state every borrowed lemma in the notation used here.

\paragraph*{How the results are checked.}
Every statement carries one of the labels \proofst, \derivst, \numst\ or \conjst.
There is no computer-assisted step in the proof of Theorem~\ref{thm:main}: every constant is a closed form.
We nevertheless recomputed all of them with outward-rounded interval arithmetic (App.~\ref{app:constants}), and we evaluated every joint of the chain, the conclusion of one lemma inserted into the hypothesis of the next, on $6062$ recorded states (App.~\ref{app:joints}).
That check found three printed rounded values whose direction was wrong in the working record; the constants derived from them hold when recomputed.
It also reproduces, as a deliberately wrong version, an earlier slip in a power of $\sigma$, which fails by a factor of $9\times10^{71}$.

\paragraph*{Plan.}
We refer to Refs.~\cite{E}, \cite{F} and~\cite{H,I,J} as Papers~E, F, H, I and~J, the names used in the code and the records.
Section~\ref{sec:main} states the setting and Theorem~\ref{thm:main}.
Section~\ref{sec:chain} shows the chain.
Sections~\ref{sec:circles}--\ref{sec:inside} prove the three new steps, and Section~\ref{sec:assembly} assembles the proof and discusses the range of $\sigma$.
Section~\ref{sec:outlook} lists what remains between Theorem~\ref{thm:main} and the conjecture~\eqref{eq:conj1}.
